
\medskip

\emph{Principaux domaines abordés}: Probabilités conditionnelles et indépendance. Variables aléatoires.

\medskip

%Lors d'une kermesse, un organisateur de jeux dispose, d'une part, d'une roue comportant quatre cases blanches et huit cases rouges et, d'autre part, d'un sac contenant cinq jetons portant les numéros 1, 2, 3, 4 et 5.
%
%Le jeu consiste à faire tourner la roue, chaque case ayant la même probabilité d'être obtenue, puis à extraire un ou deux jetons du sac selon la règle suivante :
%
%\setlength\parindent{1cm}
%\begin{itemize}
%\item[$\bullet~~$] si la case obtenue par la roue est blanche, alors le joueur extrait un jeton du sac ;
%\item[$\bullet~~$] si la case obtenue par la roue est rouge, alors le joueur extrait successivement et sans remise deux jetons du sac.
%\end{itemize}
%\setlength\parindent{0cm}
%
%Le joueur gagne si le ou les jetons tirés portent tous un numéro impair.
%
%\medskip

\begin{enumerate}
\item %Un joueur fait une partie et on note $B$ l'évènement \og la case obtenue est blanche \fg,
$R$ l'évènement \og la case obtenue est rouge\fg{} et $G$ l'évènement \og le joueur gagne la partie \fg. 
	\begin{enumerate}
		\item %Donner la valeur de la probabilité conditionnelle $P_B(G)$.
Si la case est blanche on tire 1 seul jeton ; comme il y a 3 résultats impairs sur 5 numéros on a $P_B(G) = \dfrac{3}{5} = \dfrac{6}{10} = 0,6$.
		\item %On admettra que la probabilité de tirer successivement et sans remise deux jetons impairs est égale à $0,3$. 
		
%Recopier et compléter l'arbre de probabilité suivant:
Tout d'abord on a $P(B) = \dfrac{4}{12} = \dfrac13$ et donc $P(R) = \dfrac{8}{12} = \dfrac23$.

Si la case est rouge on tire successivement deux jetons : il y a $5 \times 4 = 20$ issues différentes depuis 1--2 jusqu'à 5-4 et parmi celles-ci les issues gagnantes : 1-3 ; 1-5 ; 3-1 ; 3-5 ; 5-1 et 5-3.

On a donc $P_R(G) = \dfrac{6}{20} = \dfrac{3}{10} = 0,3$ (voir l'énoncé).

On peut donc compléter l'arbre pondéré :

\begin{center}
\pstree[treemode=R,nodesepB=3pt,levelsep=3cm]{\TR{}}
{\pstree{\TR{$B$~~} \taput{$\frac13$}}
	{\TR{$G$}\taput{0,6}
	\TR{$\overline{G}$}\tbput{0,4}
	}
\pstree{\TR{$R$~~} \tbput{$\frac23$}}
	{\TR{$G$}\taput{0,3}
	\TR{$\overline{G}$}\tbput{0,7}
	}	
}
\end{center}

	\end{enumerate}

\item
	\begin{enumerate}
		\item %Montrer que $P(G) = 0,4$.
D'après la loi des probabilités totales :
		
$P(G) = P(B \cap G) + P(R \cap G)$
	
$\bullet~~$$P(B \cap G) = P(B) \times P_B(G) = \dfrac13 \times 0,6 = \dfrac{0,6}{3} = 0,2$ ; 

$\bullet~~$$P(R \cap G) = P(R) \times P_R(G) = \dfrac23 \times 0,3 = \dfrac{0,6}{3} = 0,2$.

Donc $P(G) = 0,2 + 0,2 = 0,4$.
		\item %Un joueur gagne la partie.
		
%Quelle est la probabilité qu'il ait obtenu une case blanche en lançant la roue?
Il faut trouver $P_G(B) = \dfrac{P(G \cap B)}{P(G)} = \dfrac{P(B \cap G)}{P(G)} = \dfrac{0,2}{0,4} = \dfrac2 = 0,5$.
	\end{enumerate}
\item %Les évènements $B$ et $G$ sont-ils indépendants ? Justifier.
$\bullet~~$$P(B) \times P(G) = \dfrac13 \times 0,4 = \dfrac{0,4}{3} = \dfrac{4}{30} = \dfrac{2}{15}$ ;

$\bullet~~$$P(B \cap G) = \dfrac13 \times 0,6 = 0,2 = \dfrac{2}{10} = \dfrac15$.

Donc $P(B) \times P(G) \ne P(B \cap G)$ : les évènements ne sont pas indépendants.
\item Un même joueur fait dix parties. Les jetons tirés sont remis dans le sac après chaque partie.

On note $X$ la variable aléatoire égale au nombre de parties gagnées.
	\begin{enumerate}
		\item %Expliquer pourquoi $X$ suit une loi binomiale et préciser ses paramètres.
	Les 10 épreuves sont indépendantes et à chacune la probabilité de gagner est égale à 0,4. La variable aléatoire $X$ suit donc une loi binomiale $\mathcal{B}(10~;~0,4)$.
		\item %Calculer la probabilité, arrondie à $10^{-3}$ près, que le joueur gagne exactement trois parties sur les dix parties jouées.
On a $P(X = 3 = \binom{10}{3} \times 0,4^3 \times (1 - 0,4)^7 \approx 0,215$ au millième près.
		\item %Calculer $P (X \geqslant 4)$ arrondie à $10^{-3}$ près.
		La calculatrice donne $P(X \leqslant 3) \approx 0,3823$, donc $P(X \geqslant 4) \approx 0,6177$, soit 0,618 au millième près.
		
%Donner une interprétation du résultat obtenu.

Il y a plus de 6 chances sur 10 de gagner au moins 4 parties.
	\end{enumerate}
\item %Un joueur fait $n$ parties et on note $p_n$ la probabilité de l'évènement \og le joueur gagne au moins une partie \fg.
	\begin{enumerate}
		\item %Montrer que $p_n = 1 - 0,6^n$.
On a $p_n = P(X \geqslant 1) = 1 - P(X = 0) = 1 - 0,6^n$.
		\item %Déterminer la plus petite valeur de l'entier $n$ pour laquelle la probabilité de gagner au moins une partie est supérieure ou égale à $0,99$.
		
Il faut trouver le plus petit entier $n$ tel que $p_n \geqslant 0,99$, soit :
		
$1 - 0,6^n \geqslant 0,99 \iff 0,01 \geqslant 0,6^n$ soit par croissance de la fonction logarithme : $\ln 0,01 \geqslant n \ln 0,6 \iff \dfrac{\ln 0,01}{\ln 0,6} \leqslant n$, car $\ln 0,6 < 0$.
		
Or $\dfrac{\ln 0,01}{\ln 0,6} \approx 9,02$.
		
Il faut donc jouer au moins 10 parties pour avoir une probabilité d'en gagner une avec une probabilité d'au moins 99\,\%.
	\end{enumerate}
\end{enumerate}

\bigskip

